\section{Case Studies}
\label{sec:cases}

We develop four case studies illustrating how district count, geographic
structure, and resolution interact to shape algorithmic map outcomes.

\subsection{Washington: 98 Districts, Block-Group Resolution}
\label{sec:cases:wa}

Washington's House of Representatives has 98 districts ($k=98$) with a 2020
population of 7{,}256{,}236, yielding an ideal district population of $T^* = 74{,}043$.
Washington has 1{,}458 census tracts ($k/n = 0.067 > 0.05$), requiring
block-group resolution (4{,}874 block groups).

The geographic challenge in Washington is the Cascades Range, which divides the
state into a wet, populous western region (Puget Sound metropolitan area) and a
dry, sparse eastern region. The recursive bisection algorithm naturally partitions
the state along this east-west divide at the first level, producing a western half
with approximately 83 districts and an eastern half with 15 districts.

\textbf{Results.} Mean PP: 0.412 (above the 50-state average of 0.401 for
block-group resolution). Maximum population deviation: 0.31\%. County splits:
11 (vs.\ 14 in the enacted 2021 map). Runtime: 180 seconds (block-group resolution).

The high compactness reflects Washington's relatively uniform population
distribution in the western half: Puget Sound-area districts can be drawn as
compact near-squares around cities and suburbs without the elongated coastal
configurations that reduce compactness in eastern and southern states.

\subsection{Texas: 150 Districts, Tract Resolution}
\label{sec:cases:tx}

Texas's House of Representatives has 150 districts ($k=150$) with a 2020
population of 29{,}317{,}186, yielding $T^* = 195{,}448$. Texas has 5{,}265
census tracts ($k/n = 0.028$), below the threshold for block-group resolution.

The geographic challenge in Texas is extreme: four major metropolitan areas
(Houston, Dallas-Fort Worth, San Antonio, Austin) contain over 70\% of the
population in roughly 15\% of the state's land area. The algorithm must draw
a large number of compact urban districts in dense areas while also covering
the vast, sparsely populated West Texas and Panhandle regions with geographically
large but compact districts.

\textbf{Results.} Mean PP: 0.374 (below the 50-state average, reflecting Texas's
challenging urban-rural geography). Maximum population deviation: 0.44\%.
County splits: 48 (vs.\ 62 in the enacted 2021 map). Runtime: 95 seconds.

Texas illustrates a general pattern: states with extreme population concentration
in a few metropolitan areas achieve lower average compactness because urban
districts must be carved into compact but geometrically constrained shapes.
The 23\% reduction in county splits (48 vs.\ 62) is consistent with the
algorithm's preference for geographic coherence over political subdivision
boundaries.

\subsection{New Hampshire: 400 Districts, Block-Group Resolution}
\label{sec:cases:nh}

New Hampshire's House of Representatives is the largest state legislative
chamber in the U.S.\ and one of the largest deliberative bodies in the world,
with 400 members. With a 2020 population of 1{,}377{,}529, the ideal district
population is $T^* = 3{,}444$---roughly the size of a small town.
New Hampshire has only 321 census tracts, making redistricting at tract
resolution impossible ($k/n = 1.25 > 1$). We use block-group resolution:
New Hampshire has 698 block groups, giving a ratio of $k/n_{\rm bg} = 0.57$.
This ratio is still above the optimal 20-units-per-district threshold, but
block-group redistricting is feasible and produces valid non-empty districts.

\textbf{Results.} Mean PP: 0.388. Maximum population deviation: 0.48\%.
County splits: 31 (New Hampshire has 10 counties, so this represents a
moderate level of splitting relative to the large number of districts).
Runtime: 420 seconds.

New Hampshire illustrates the hardest case for the pipeline. With $k/n_{\rm bg} = 0.57$,
the average block group contains only $1.75$ districts---meaning many districts
must consist of a single block group or split block groups. The pipeline handles
this correctly via population-weighted sub-unit assignment when a block group
is the only unit available to complete a district.

\subsection{Nebraska: Unicameral Legislature, 49 Districts}
\label{sec:cases:ne}

Nebraska is the only unicameral state legislature in the U.S. The Nebraska
Legislature has 49 members (called senators), with a 2020 population of
1{,}961{,}504 and ideal district population $T^* = 40{,}031$.
Nebraska has 570 census tracts ($k/n = 0.086 > 0.05$), requiring block-group
resolution (1{,}231 block groups, $k/n_{\rm bg} = 0.040$).

\textbf{Results.} Mean PP: 0.421 (highest among the block-group-resolution states).
Maximum population deviation: 0.38\%. County splits: 8. Runtime: 165 seconds.

Nebraska's high compactness reflects its geography: the state is roughly
rectangular, population is distributed along the Platte River corridor from
Omaha to Lincoln westward, and the Sandhills region occupies a large but very
sparsely populated central area. The algorithm naturally groups Omaha-area
block groups into compact eastern districts and distributes the few population
centres of western Nebraska into geographically large but compactly bounded
districts.

Because Nebraska has no upper chamber, the NestSection algorithm (F.2) does
not apply. This simplifies the problem relative to other states: there is
no constraint to maintain senate-house nesting.
